\section[Notebook 1973-6: commutators of elementary matrices; a Sylow reduction]{Notebook 1973-6: commutators of elementary matrices, and a Sylow reduction\protect\verifie{Notebook1973_6}}\label{sec:quillen-1973-6}

\emph{Context.} \carnet{1973}{6} holds Quillen's notes of Serre's course on trees, in the autumn of 1973, followed by his own pages on shift equivalence. Two short arguments of the course are checked here.

\emph{The page, p.~30.} To show that $SL_3(\mathbf Z)$ has Serre's property (FA), the page writes two commutator relations between the elementary matrices $\bar e_{ij} = 1 + e_{ij}$, for $i, j, k$ distinct: $(\bar e_{ij}, \bar e_{jk}) = \bar e_{ik}$ and $(\bar e_{ij}, \bar e_{kj}) = \bar e_{ik}^{-1}$, with $(a, b) = aba^{-1}b^{-1}$. It then places the six matrices on a hexagon, in the order $\bar e_{12}, \bar e_{13}, \bar e_{23}, \bar e_{21}, \bar e_{31}, \bar e_{32}$, and notes that each commutes with its neighbours and is the commutator of its two neighbours. The first relation is the usual Steinberg relation. The second cannot hold as written: $\bar e_{ij}$ and $\bar e_{kj}$ lie in the same column, $e_{ij}e_{kj} = e_{kj}e_{ij} = 0$, and they commute. The relation that holds is $(\bar e_{ij}, \bar e_{ki}) = \bar e_{kj}^{-1}$; the hexagon needs only the first one.

\emph{The page, p.~45.} In the proof of Brown's theorem that an integer dividing the length of every orbit of a finite group divides the Euler characteristic, the page reduces to a Sylow $p$-subgroup $H$: if $p^k$ divides $[G : G_x]$, it divides $[H : H_x]$, because $[G:G_x]$ divides $[G : H_x] = [G : H][H : H_x]$ and $[G : H]$ is prime to $p$.

\begin{encadre}
\begin{theoreme}[After Quillen, notebook 1973-6, p.~30]\label{thm:quillen-1973-6-a}
Let $R$ be a commutative ring, $i, j, k$ distinct indices, $a, b \in R$, and $t_{ij}(a) = 1 + a e_{ij}$. Then $(t_{ij}(a), t_{jk}(b)) = t_{ik}(ab)$; $t_{ij}(a)$ commutes with $t_{kj}(b)$ and with $t_{ik}(b)$, so that $(t_{ij}(a), t_{kj}(b)) = 1$, which is not $t_{ik}(-1)$ if $R \ne 0$; and $(t_{ij}(a), t_{ki}(b)) = t_{kj}(-ba)$. For $3 \times 3$ matrices, each matrix of the hexagon commutes with the next one and is the commutator, in one order or the other, of its two neighbours.
\end{theoreme}
\end{encadre}

\begin{proof}
Put $x = ae_{ij}$ and $y = be_{jk}$. Then $x^2 = y^2 = yx = 0$ and $xy = ab\,e_{ik}$, $xyx = 0$; expanding, $(1+x)(1+y)(1-x)(1-y) = 1 + xy$. For $y = be_{ki}$ instead, $xy = 0$, $yx = ba\,e_{kj}$ and $yxy = 0$, and the same expansion gives $1 - yx$. Two elementary matrices in the same row or the same column have products $0$ in both orders, so they commute. The hexagon is the six instances of the first relation: $\bar e_{12} = (\bar e_{13}, \bar e_{32})$, $\bar e_{13} = (\bar e_{12}, \bar e_{23})$, $\bar e_{23} = (\bar e_{21}, \bar e_{13})$, and so on, and consecutive matrices share a row or a column.
\end{proof}

\begin{encadre}
\begin{theoreme}[Quillen, notebook 1973-6, p.~45]\label{thm:quillen-1973-6-b}
Let $H$ be a Sylow $p$-subgroup of a finite group $G$ and $K$ a subgroup. If $p^k$ divides $[G : K]$, then $p^k$ divides $[H : H \cap K]$.
\end{theoreme}
\end{encadre}

\begin{proof}
$[G : K]$ divides $[G : H \cap K] = [H : H \cap K]\,[G : H]$, and $[G:H]$ is prime to $p$.
\end{proof}

\emph{Verification in Lean.} The file \href{\leanblobmain/Notebook1973_6.lean}{\texttt{Notebook1973\_6.\allowbreak{}lean}} proves Theorem~\ref{thm:quillen-1973-6-a} with mathlib's \texttt{Matrix.transvection}, as \texttt{transvection\_commutator}, \texttt{transvection\_commute\_col}, \texttt{transvection\_commute\_row}, \texttt{transvection\_commutator\_col}, \texttt{not\_second\_relation}, \texttt{transvection\_commutator'} and, for the hexagon, \texttt{hexagon\_commute} and \texttt{hexagon\_commutator}; the commutators are written out as products, the inverse of $t_{ij}(a)$ being $t_{ij}(-a)$. The first relation was already in the library for units over any ring (Section~\ref{sec:quillen-1968-15}). Theorem~\ref{thm:quillen-1973-6-b} is \texttt{pow\_dvd\_relIndex\_of\_sylow}. The property (FA), the trees and Brown's theorem are not formalised. The axioms used are the three usual axioms of mathlib.

\litterature{Both statements are classical: the commutator relations are Steinberg's relations for $SL_n$, and the hexagon argument is in Serre's \emph{Arbres, amalgames, $SL_2$} (I, \S 6); the Sylow reduction is the first step of Brown's proof (Invent.\ Math.\ 27, 1974). The point of the section is the second relation of p.~30, which fails as transcribed and is replaced by the one that holds. The edition does not present anything here as new.}
